\section{Deferred Statements and Proofs}
\label{app:proofs}

\begin{lemma}[Gauge]
\label{lem:gauge}
Let $\mathcal{L}_{\mathrm{rel}}$ depend on the student only through
$\{\cos(s_i,s_j)\}_{i,j}$; its invariance group contains
$G_{\mathrm{rel}} = \{s\mapsto cRs : c>0,\ R\in O(d_S)\}$. For
$\mathcal{L}_{\mathrm{anc}} = \frac{1}{|B|}\sum_{i\in B}\big(1-\cos(W_a s_i,t_i)\big)$
with $W_a\in\R^{d_T\times d_S}$ free, the substitution $s\mapsto As$,
$W_a\mapsto W_a A^{-1}$ leaves the term exactly invariant for every $A\in GL(d_S)$. Hence
$G_{\mathrm{anc}} = GL(d_S) \supsetneq G_{\mathrm{rel}}$, and
$\min_{W_a}\mathcal{L}_{\mathrm{anc}}$ is a function of the student space only through its
$GL(d_S)$-orbit.
\end{lemma}

\begin{proof}
For any invertible $A$, $W_a A^{-1}(As_i) = W_a s_i$, so every summand is unchanged; the
map $W_a\mapsto W_aA^{-1}$ is a bijection of the parameter space, so the minimised value is
unchanged as well. Taking $A=cR$ recovers $G_{\mathrm{rel}}\subseteq G_{\mathrm{anc}}$, and
any $A$ with distinct singular values lies in $G_{\mathrm{anc}}\setminus G_{\mathrm{rel}}$,
which gives strict containment.
\end{proof}

\begin{proof}[Proof of Proposition~\ref{prop:levels}]
(a) is the shift invariance of the softmax: replacing $v_i$ by $v_i+a\mathbf{1}$ multiplies
numerator and denominator of \eqref{eq:readout} by $e^{a/\tau}$.

(b) Exact satisfaction $q_i^{\tau_r}=p_{i,r}$ holds iff
$v_i/\tau_r = \log p_{i,r} + c_r\mathbf{1}$ for some $c_r\in\R$, i.e.\
$v_i = \tau_r\log p_{i,r} + \tau_r c_r\mathbf{1}$. Equating two such expressions for
$r\neq r'$ gives \eqref{eq:consistency}. The first constraint already pins $v_i$ up to one
scalar; each further $r$ imposes $C-1$ additional conditions on that single remaining
degree of freedom, which for generic targets admit no solution. The minimiser of
$\sum_r\omega_r\KL(p_{i,r}\Vert q_i^{\tau_r})$ is therefore a compromise, and since the
residuals it balances are the vectors $\tau_r\log p_{i,r}$ themselves, its location depends
on the absolute values of $\log p_{i,r}$ and not only on the shift-equivalence class of
$v_i$.
\end{proof}

\section{Implementation Details}
\label{app:impl}

This appendix fixes the constants left symbolic in the main text. Nothing here is required
by the analysis of Section~\ref{sec:analysis}; the choices are reported so that
Table~\ref{tab:main} is reproducible and so that Condition~\ref{cond:cover} can be checked
rather than assumed.

\begin{algorithm}[h]
\caption{\method training}
\label{alg:method}
\begin{algorithmic}[1]
\REQUIRE Corpus $\mathcal{D}$, teacher $T$, student $S_\theta$, graph size $k$, scales
$\Rset=\{0\}\cup\mathcal{R}$, temperatures $\{\tau_r\}$, weights $\{\omega_r\}$, pool size
$P$, candidate size $C$
\STATE Cache $t_i = T(x_i)$ for all $i$; discard $T$
\STATE Build $W$ by \eqref{eq:graph}; form $P$ and $Q$ as in Section~\ref{sec:setup}
\STATE Precompute pools $\mathcal{P}_i$ and targets $p_{i,r}$, $r\in\mathcal{R}$, from one walk snapshotted at each $r$; record residuals $\varepsilon_{i,r}$
\STATE Compute the floor \eqref{eq:floor} and the build-time diagnostics \eqref{eq:diag1}
\FOR{each epoch $e$}
  \STATE Redraw $C_i^{(e)}\sim\mu_i$ for every anchor $i$ \COMMENT{Condition~\ref{cond:cover}}
  \FOR{each mini-batch $B$}
    \STATE Encode anchors and candidates with $S_\theta$; form the shared column set $\mathcal{C}_B$ and scatter targets onto it
    \STATE Evaluate $p_{i,0}$ on $\mathcal{C}_B$ from the cached teacher embeddings
    \STATE Compute $q_i^{\tau_r}$ for every $r\in\Rset$ and update $\theta$ on \eqref{eq:objective}
  \ENDFOR
\ENDFOR
\RETURN $S_\theta$
\end{algorithmic}
\end{algorithm}

\subsection{Hyperparameters}

\begin{table}[h]
\centering
\small
\begin{tabular}{llc}
\toprule
Symbol & Meaning & Value \\
\midrule
$k$ & mutual $k$NN graph degree cap & $50$ \\
$\tau_g$ & graph kernel temperature, \eqref{eq:graph} & $0.05$ \\
$\tau_0^{\mathrm{tgt}}$ & teacher temperature of the $r{=}0$ target, \eqref{eq:family} & $0.10$ \\
$\mathcal{R}$ & diffusion scales & $\{1,2,4\}$ \\
$(\tau_1,\tau_2,\tau_4)$ & readout temperatures, diffusion scales & $(0.05,\,0.07,\,0.10)$ \\
$\tau_0$ & readout temperature, zero-step scale & $0.10$ \\
$(\omega_0,\omega_1,\omega_2,\omega_4)$ & scale weights, normalised & $(0.500,\,0.286,\,0.143,\,0.071)$ \\
$P$ & retained pool size per anchor & $128$ \\
$C$ & candidate set size per anchor & $64$ \\
--- & optimiser / LR / batch / epochs & AdamW / $2{\times}10^{-5}$ / $16$ / $5$ \\
\bottomrule
\end{tabular}
\end{table}

The diffusion readout temperatures increase with $r$, so that broader targets are read out
more smoothly. The zero-step readout temperature is set independently rather than by
extending that trend: $q_i^{\tau_0}$ ranges over the entire candidate set rather than a
$k$NN neighbourhood, where inheriting $\tau_g$ would place nearly all mass on the nearest
column, and it is paired with a target whose own temperature is $\tau_0^{\mathrm{tgt}}$.
Proposition~\ref{prop:levels}(b) requires only that at least two readout temperatures
differ, which the bank satisfies. The graph temperature $\tau_g$ is selected by the
sharpness diagnostic of \eqref{eq:diag1} rather than by a downstream sweep: teacher cosines
within an anchor's top-$k$ neighbourhood occupy a narrow band, so a $\tau_g$ large relative
to that band flattens every row of $P$ towards uniform on its support, which sharpness
detects offline.

The weights follow from normalising the raw diffusion weights $(1,\,\tfrac12,\,\tfrac14)$
to sum to one and then admitting $r=0$ with unit relative weight, so the zero-step scale
carries half the total. We record this explicitly because it is easy to introduce a scale
whose effective weight is far from what a list of raw coefficients suggests, and because
$\omega$ enters both the objective \eqref{eq:objective} and the floor \eqref{eq:floor} ---
a floor computed under different weights than the loss is not comparable to it.

\subsection{The candidate sampler}

Condition~\ref{cond:cover} is realised by a fixed quota allocation of the $C=64$ slots:
\begin{itemize}
    \item $32$ diffusion slots, split evenly across $\mathcal{R}$ (not drawn from the
    mixture $\bar p_i$), so each scale is guaranteed scored columns from its own support;
    \item $24$ teacher hard negatives, drawn from the $200$ nearest teacher neighbours that
    lie outside the graph pool;
    \item $8$ uniform random negatives.
\end{itemize}
Within a scale's quota, the top $2$ columns by diffusion mass are taken deterministically
and the remainder are drawn without replacement by Gumbel top-$k$ over
$\log (P_r)_{ij}$, which samples proportionally to diffusion mass while keeping the draw
vectorised. Sets are redrawn once per epoch, for a reason independent of
Condition~\ref{cond:cover}: a frozen candidate set defines a fixed finite collection of
comparisons, so once the student orders them correctly the gradient vanishes even though
$p_{i,r}$ still describes unlearned structure. If mutuality isolates an anchor --- $2.0\%$ of
anchors on our corpus --- its pool falls back to the teacher's directed top-$k$ neighbours;
these anchors still receive a well-defined $r=0$ target, which is one practical consequence
of Proposition~\ref{prop:support}'s remedy.

\subsection{Corpus preparation and targets}

Exact-duplicate texts are removed before graph construction ($13.5$k anchors remain): they
have $\cos=1$ by construction, so their logits sit at the ceiling with vanishing gradient
while still absorbing teacher mass from candidates that could still move.
Diffusion targets are precomputed once by a single sparse walk snapshotted at each
$r\in\mathcal{R}$, retaining the top $P=128$ entries per row and recording the truncation
residual $\varepsilon_{i,r}$ of \eqref{eq:target}; the $r=0$ target is evaluated on the fly
from cached teacher vectors, so it is exact on $\mathcal{C}_B$ and requires no storage. All
targets are renormalised on the scored column set after masking the anchor's own index.

\subsection{Build-time diagnostics}
\label{app:diag}

% Build-time diagnostics, Qwen3-Embedding-4B teacher, 13.5k deduplicated anchors.
% All quantities computed from the cached teacher artifact before any gradient step.
\begin{table}[t]
\centering
\caption{Build-time diagnostics of the scale family \eqref{eq:family} for the
Qwen3-Embedding-4B teacher ($k=50$, $\tau_g=0.05$), averaged over anchors and evaluated on
the retained pool. Sharpness is the first quantity in \eqref{eq:diag1}, in nats; a value near $0$ would indicate
a binary neighbour label. Residual is $\varepsilon_{i,r}$ of \eqref{eq:target}. The
zero-step scale is dense on $C_i$ by construction and is reported separately in the text.}
\label{tab:diag}
\small
\begin{tabular}{lcccc}
\toprule
Scale & $|\supp(p_{i,r})|$ & Sharpness \eqref{eq:diag1} & Top-1 mass & Residual $\varepsilon_{i,r}$ \\
\midrule
$r=1$ & \phantom{0}20.5 & 0.649 & 0.401 & $1.4\times10^{-17}$ \\
$r=2$ & 101.8 & 1.719 & 0.333 & $2.4\times10^{-2}$ \\
$r=4$ & 124.9 & 1.528 & 0.257 & $8.7\times10^{-2}$ \\
\midrule
\multicolumn{5}{l}{\emph{Scale distinctness}: $\KL(p_{i,1}\Vert p_{i,2})=0.173$, \ $\KL(p_{i,2}\Vert p_{i,4})=0.127$} \\
\multicolumn{5}{l}{\emph{Irreducible floor} \eqref{eq:floor}: $\JS_\omega = 0.100$ (mean), $0.191$ (90th percentile)} \\
\multicolumn{5}{l}{\emph{Graph}: mean degree $20.5$; $2.0\%$ of anchors isolated by mutuality and restored by fallback} \\
\bottomrule
\end{tabular}
\end{table}


\subsection{Logged quantities}

Per epoch we log: total loss and per-scale KL; the excess $\mathcal{L}-\JS_\omega$; student
entropy relative to $\log|\mathcal{C}_B|$ per scale; student mass on the teacher's top-$8$
candidates; mean and $90$th-percentile truncation residual; and the number of steps skipped
for non-finite gradients. The last is reported because a silently skipped step is
indistinguishable from a completed one in the loss curve. Checkpoint selection uses a
held-out validation metric, not the training objective: by Corollary~\ref{cor:floor} the
training loss contains a large anchor-dependent constant, so its across-epoch differences
are not a reliable ordering of model quality.
